% !TEX program = pdflatex
\documentclass[handout,aspectratio=169,12pt]{beamer}

% =========================
% PAQUETES
% =========================
\usepackage[spanish,es-noshorthands]{babel}
\usepackage[T1]{fontenc}
\usepackage[utf8]{inputenc}
\usepackage{lmodern}
\usepackage{amsmath,amssymb,amsfonts}
\usepackage{graphicx}
\usepackage{xcolor}
\usepackage{booktabs}
\usepackage{tikz}
\usepackage{array}

% =========================
% TEMA Y COLORES
% =========================
\usetheme{Madrid}
\usecolortheme{default}
\usefonttheme[onlymath]{serif}

\definecolor{USTBlue}{RGB}{0,70,127}
\definecolor{USTBlueDark}{RGB}{0,45,85}
\definecolor{USTGreen}{RGB}{27,94,32}
\definecolor{USTBrown}{RGB}{93,64,55}
\definecolor{USTSoft}{RGB}{248,249,250}
\definecolor{USTText}{RGB}{35,40,45}
\definecolor{USTAccent}{RGB}{21,101,192}

\setbeamercolor{title}{fg=white,bg=USTBlue}
\setbeamercolor{frametitle}{fg=white,bg=USTBlue}
\setbeamercolor{structure}{fg=USTBlue}
\setbeamercolor{normal text}{fg=USTText,bg=white}
\setbeamercolor{block title}{fg=white,bg=USTGreen}
\setbeamercolor{block body}{fg=black,bg=USTSoft}
\setbeamercolor{block title example}{fg=white,bg=USTAccent}
\setbeamercolor{block body example}{fg=black,bg=blue!4}
\setbeamercolor{block title alerted}{fg=white,bg=USTBrown}
\setbeamercolor{block body alerted}{fg=black,bg=brown!8}

\setbeamertemplate{navigation symbols}{}

\setbeamertemplate{footline}{
	\leavevmode%
	\hbox{%
		\begin{beamercolorbox}[wd=.77\paperwidth,ht=2.6ex,dp=1ex,leftskip=0.4cm]{author in head/foot}%
			\color{white}\ifx\insertshorttitle\empty\else\insertshorttitle\fi
		\end{beamercolorbox}%
		\begin{beamercolorbox}[wd=.23\paperwidth,ht=2.6ex,dp=1ex,rightskip=0.35cm plus1fil]{date in head/foot}%
			\color{white}\insertframenumber/\inserttotalframenumber
	\end{beamercolorbox}}%
}
\setbeamercolor{author in head/foot}{bg=USTBlueDark,fg=white}
\setbeamercolor{date in head/foot}{bg=USTBlueDark,fg=white}

% =========================
% DATOS
% =========================
\title[Fundamentos de Matemática Básica]{Fundamentos de Matemática Básica}
\subtitle{Unidad I \\ Clase 6: Expresiones Algebraicas}
\author{Profesor Luis Tulio González Alcaino \\ \small Magíster en Matemática}
\institute{Universidad Santo Tomás \\ Departamento Ciencias Básicas -- Talca}
\date{Semestre I -- 2026}

% =========================
% COMANDOS
% =========================
\newcommand{\IdeaClave}[1]{
	\begin{alertblock}{Idea clave}
		#1
	\end{alertblock}
}

\newcommand{\Agronomia}[1]{
	\begin{exampleblock}{Contexto agronómico}
		#1
	\end{exampleblock}
}

% =========================
% DOCUMENTO
% =========================
\begin{document}
	
	% ---------------------------------
	\begin{frame}
		\titlepage
	\end{frame}
	
	% ---------------------------------
\begin{frame}{Ruta de la clase}
	
	\begin{block}{¿Qué aprenderemos hoy?}
		\begin{itemize}[<+->]
			\item Comprender qué es un término algebraico
			\item Identificar expresiones algebraicas
			\item Reconocer términos semejantes
			\item Simplificar expresiones
		\end{itemize}
	\end{block}
	
	\pause
	
	\begin{exampleblock}{¿Dónde lo aplicamos?}
		\begin{itemize}[<+->]
			\item Modelación de riego
			\item Fertilización
			\item Análisis de datos agrícolas
		\end{itemize}
	\end{exampleblock}
	
\end{frame}

	\section{Motivación}
	% ---------------------------------
	\begin{frame}{Motivación}
		\begin{block}{Situación inicial}
			En un cultivo de lechugas:
			\begin{itemize}
				\item Se aplican \(2x\) litros de agua por planta en condiciones normales.
				\pause
				\item Se agregan \(3x\) litros en días calurosos.
				\pause
				\item Se pierde \(x\) litro por evaporación.
			\end{itemize}
		\end{block}
		
		\pause
		\textbf{Pregunta:} ¿Cuál es el aporte real de agua por planta?
		
		\pause
		$
		2x+3x-x
		$
		
		\pause
		$
		5x-x=4x
		$
		
		\pause
		\IdeaClave{El aporte real de agua por planta es \(4x\) litros.}
	\end{frame}
	
	% ---------------------------------
	\begin{frame}{Ejemplo resuelto más desafiante}
		\Agronomia{
			En una parcela experimental de tomates:
			\begin{itemize}
				\item en la mañana se aplican \(4x\) litros de solución nutritiva,
				\pause
				\item al mediodía se agregan \(3x\) litros,
				\pause
				\item por escurrimiento se pierden \(2x\) litros,
				\pause
				\item y por corrección técnica se añaden \(x\) litros extra.
			\end{itemize}
		}
		
		\pause
		\textbf{Modelo algebraico: } $4x+3x-2x+x$
		
		\medskip
		
		\pause
		
		\textbf{Simplificación:}
	
	\medskip 
	
	\qquad $4x+3x+x-2x=8x-2x=6x$
		
		\pause
		\IdeaClave{El aporte neto total es \(6x\) litros de solución nutritiva.}
	\end{frame}
	
	\section{Término algebraico}
	% ---------------------------------
	\begin{frame}{¿Qué es un término algebraico?}
		\begin{block}{Definición}
			Un \textbf{término algebraico} es una expresión formada por:
			\begin{itemize}
				\item un número, llamado \textbf{coeficiente},
				\pause
				\item una o más \textbf{variables},
				\pause
				\item un \textbf{exponente}, si corresponde.
			\end{itemize}
		\end{block}
		
		\pause
		\Agronomia{
			El término \(5x\) puede representar 5 litros de fertilizante por cada hectárea \(x\).
		}
	\end{frame}
	
	% ---------------------------------
	\begin{frame}{Ejemplos de términos algebraicos en agronomía}
		\begin{columns}[T]
			\begin{column}{0.48\textwidth}
				\begin{block}{Ejemplo 1}
					\[
					7x
					\]
					\pause
					Representa 7 kg de semilla por lote \(x\).
				\end{block}
				
				\pause
				\begin{block}{Ejemplo 2}
					\[
					3y^2
					\]
					\pause
					Puede representar una cantidad que depende del cuadrado de la superficie sembrada.
				\end{block}
			\end{column}
			
			\begin{column}{0.48\textwidth}
				\begin{exampleblock}{Ejemplo 3resuelto más complejo}
					\[
					-4ab^2
					\]
					\pause
					\begin{itemize}
						\item Coeficiente: \(-4\)
						\pause
						\item Variables: \(a\) y \(b\)
						\pause
						\item Exponentes: \(a^1\), \(b^2\)
					\end{itemize}
				\end{exampleblock}
			\end{column}
		\end{columns}
	\end{frame}
	
	\section{Partes de un término algebraico}
	% ---------------------------------
	\begin{frame}{Partes de un término algebraico}
		Consideremos el término:
		\[
		7x^2
		\]
		
		\pause
		\begin{itemize}
			\item \textbf{Coeficiente:} \(7\)
			\pause
			\item \textbf{Variable:} \(x\)
			\pause
			\item \textbf{Exponente:} \(2\)
		\end{itemize}
		
		\pause
		\IdeaClave{Si una variable no muestra exponente, se entiende que su exponente es \(1\).}
	\end{frame}
	
	% ---------------------------------
	\begin{frame}{Ejemplo resuelto: identificar partes}
		\Agronomia{
			En un registro de producción aparece el término:\quad $12m^3n$
			
			\smallskip 
			
			donde \(m\) representa la dosis de materia orgánica y \(n\) el número de sectores tratados.
		}
		
		\pause
		\textbf{Identifiquemos sus partes:}
		\begin{itemize}
			\item Coeficiente: \(12\)
			\pause
			\item Variables: \(m\) y \(n\)
			\pause
			\item Exponente de \(m\): \(3\)
			\pause
			\item Exponente de \(n\): \(1\)
		\end{itemize}
		
		\pause
		\IdeaClave{Un mismo término puede tener más de una variable, y cada una puede tener distinto exponente.}
	\end{frame}
	
	\section{Expresión algebraica}
	% ---------------------------------
	\begin{frame}{Expresión algebraica}
		\begin{block}{Definición}
			Una \textbf{expresión algebraica} es una combinación de términos algebraicos mediante operaciones.
		\end{block}
		
		\pause
		\[
		2x + 3x - 4
		\]
		
		\pause
		\Agronomia{
			\begin{itemize}
				\item Agua base: \(2x\)
				\pause
				\item Agua adicional: \(3x\)
				\pause
				\item Pérdida fija por evaporación: \(4\)
			\end{itemize}
		}
	\end{frame}
	
	% ---------------------------------
	\begin{frame}{Ejemplo resuelto de expresión algebraica}
		\Agronomia{
			En un cultivo de frutillas:
			\begin{itemize}
				\item se aplican \(5x\) litros de riego por goteo,
				\pause
				\item se agregan \(2x\) litros por refuerzo,
				\pause
				\item se descuentan \(3\) litros por pérdida fija,
				\pause
				\item y se suman \(x\) litros en una corrección final.
			\end{itemize}
		}
		
		\pause
		\textbf{Expresión:}\quad $5x+2x-3+x$
		
		\pause
		\textbf{Agrupando términos semejantes:}\quad $5x+2x+x-3=8x-3$
		
		\pause
		\IdeaClave{La expresión simplificada es \(8x-3\).}
	\end{frame}
	
	\section{Términos semejantes}
	% ---------------------------------
	\begin{frame}{Términos semejantes}
		\begin{block}{Definición}
			Dos o más términos son \textbf{semejantes} si tienen la misma variable elevada al mismo exponente.
		\end{block}
		
		\pause
		\[
		2x+3x=5x
		\]
		
		\pause
		\Agronomia{
			Fertilizante aplicado en la mañana: \(2x\) \\
			Fertilizante aplicado en la tarde: \(3x\) \\
			Total: \(5x\)
		}
	\end{frame}
	
	% ---------------------------------
	\begin{frame}{Ejemplo resuelto: reconocer términos semejantes}
		\Agronomia{
			Analicemos los términos de una expresión asociada al manejo de un invernadero:
			\[
			4x^2,\quad -3x^2,\quad 5x,\quad 2y,\quad 7
			\]
		}
		
		\pause
		\textbf{Observación:}
		\begin{itemize}
			\item \(4x^2\) y \(-3x^2\) son semejantes.
			\pause
			\item \(5x\) no es semejante a \(4x^2\) porque cambia el exponente.
			\pause
			\item \(2y\) no es semejante a \(5x\) porque cambia la variable.
			\pause
			\item \(7\) es un término constante.
		\end{itemize}
		
		\pause
		\textbf{Si simplificamos solo los semejantes:}\quad $4x^2-3x^2=x^2$
		
		\pause
		\IdeaClave{Solo se suman o restan directamente los términos semejantes.}
	\end{frame}
	
	% ---------------------------------
	\begin{frame}{¿Cuándo no son términos semejantes?}
		\[
		2x,\quad 3x^2,\quad 4y,\quad 5
		\]
		
		\pause
		\begin{itemize}
			\item \(2x\) y \(3x^2\) no son semejantes porque cambia el exponente.
			\pause
			\item \(2x\) y \(4y\) no son semejantes porque cambia la variable.
			\pause
			\item \(5\) es un término constante.
		\end{itemize}
		
		\pause
		\IdeaClave{No basta con que “se parezcan”; deben tener exactamente la misma parte literal y el mismo exponente.}
	\end{frame}
	
	\section{Actividades}
	% ---------------------------------
	\begin{frame}{Actividad 1}
		\textbf{Simplifique:}
		\[
		4x+6x
		\]
		
		\pause
		\textbf{Respuesta:}
		\[
		10x
		\]
		
		\pause
		\Agronomia{Interpretación: total de nutrientes aplicados en un cultivo.}
	\end{frame}
	
	% ---------------------------------
	\begin{frame}{Actividad 1 ampliada}
		\Agronomia{
			En un huerto se aplican \(4x\) kg de compost en una zona y \(6x\) kg en otra zona equivalente. Además, luego se incorporan \(2x\) kg extra por corrección nutricional.
		}
		
		\pause
		\textbf{Expresión:}
		\[
		4x+6x+2x
		\]
		
		\pause
		\[
		10x+2x=12x
		\]
		
		\pause
		\IdeaClave{La cantidad total de compost aplicada es \(12x\) kg.}
	\end{frame}
	
	% ---------------------------------
	\begin{frame}{Actividad 2}
		\textbf{Simplifique:}
		\[
		7x-2x+3
		\]
		
		\pause
		\textbf{Respuesta:}
		\[
		5x+3
		\]
		
		\pause
		\Agronomia{Interpretación: crecimiento neto más un factor constante.}
	\end{frame}
	
	% ---------------------------------
	\begin{frame}{Actividad 2 ampliada}
		\Agronomia{
			Una plantación tiene una ganancia de \(9x\) plantas por trasplante, una pérdida de \(4x\) por estrés hídrico y un aumento fijo de 6 plantas provenientes de reposición manual.
		}
		
		\pause
		\textbf{Modelo algebraico:}
		\[
		9x-4x+6
		\]
		
		\pause
		\[
		5x+6
		\]
		
		\pause
		\IdeaClave{La expresión simplificada es \(5x+6\).}
	\end{frame}
	
	% ---------------------------------
	\begin{frame}{Actividad 3}
		\textbf{Simplifique:}
		\[
		8x+5x-2x
		\]
		
		\pause
		\textbf{Respuesta:}
		\[
		11x
		\]
		
		\pause
		\Agronomia{Interpretación: cantidad total de agua o fertilizante aplicada a una zona de cultivo.}
	\end{frame}
	
	% ---------------------------------
	\begin{frame}{Actividad 3 ampliada}
		\Agronomia{
			En un ensayo agronómico se registran tres aplicaciones de bioestimulante:
			\[
			8x,\quad 5x,\quad -2x,\quad +3x
			\]
			donde el último término corresponde a un refuerzo extraordinario.
		}
		
		\pause
		\textbf{Simplificación:}
		\[
		\begin{array}{c}
		  8x+5x-2x+3x \\
		  13x-2x+3x \\
		  11x+3x=14x
		\end{array}
		\]
		\IdeaClave{El total neto aplicado es \(14x\).}
	\end{frame}
	
	% ---------------------------------
	\begin{frame}{Actividad 4 evaluación}
		Si \(x=5\), ¿cuál es el valor de la expresión?
		\[
		3x+2
		\]
		
		\pause
		\[
		3(5)+2
		\]
		
		\pause
		\[
		15+2=17
		\]
	\end{frame}
	
	% ---------------------------------
	\begin{frame}{Actividad 4 ampliada}
		\Agronomia{
			Si \(x=4\), determine el valor de la expresión que modela el consumo de agua en un microtúnel:
			\[
			2x+3x+7
			\]
		}
		
		\pause
		\textbf{Sustituyendo \(x=4\):}
		\[
		2(4)+3(4)+7
		\]
		
		\pause
		\[
		8+12+7 = 27
		\]
		
		\pause
		\IdeaClave{El valor de la expresión es \(27\).}
	\end{frame}
	
	\section{Modelación}
	% ---------------------------------
	\begin{frame}{Actividad de modelación}
		\begin{block}{Problema}
			Un cultivo tiene:
			\begin{itemize}
				\item \(4x\) plantas iniciales,
				\pause
				\item se agregan \(2x\),
				\pause
				\item se pierden \(x\).
			\end{itemize}
		\end{block}
		
		\pause
		\[
		4x+2x-x
		\]
		
		\pause
		\[
		6x-x=5x
		\]
		
		\pause
		\IdeaClave{El total final se modela por \(5x\).}
	\end{frame}
	
	% ---------------------------------
	\begin{frame}{Actividad de modelación ampliada}
		\Agronomia{
			En un vivero:
			\begin{itemize}
				\item se preparan \(6x\) bandejas de almácigo,
				\pause
				\item se suman \(3x\) bandejas nuevas,
				\pause
				\item se descartan \(2x\) por baja germinación,
				\pause
				\item y luego se recuperan \(x\) bandejas mediante resiembra.
			\end{itemize}
		}
		
		\pause
		\textbf{Modelo:}
		\[
		\begin{array}{c}
		6x+3x-2x+x \\
		6x+3x+x-2x = 8x \\
		\end{array}
		\]
		\pause
		\IdeaClave{La cantidad final de bandejas operativas es \(8x\).}
	\end{frame}
	
	\section{Ejercicios}
	% ---------------------------------
	\begin{frame}{Ejercicios}
		\begin{enumerate}[<+->]
			\item Simplifique: \(8x+2x-5x\)
			\item Simplifique: \(3xy+4xy\)
			\item Evalúe \(5x+1\) si \(x=3\)
			\item Detecte patrón: \(x,3x,9x,27x\)
			\item Detecte inconsistencia: \(2x,4x,6x,11x\)
		\end{enumerate}
	\end{frame}
	
	% ---------------------------------
	\begin{frame}{Ejercicios resueltos y ampliados}
		\small
		\begin{enumerate}
			\item \textbf{Simplifique:} \(8x+2x-5x\)
			
			\pause
			\[
			10x-5x=5x
			\]
			
			\pause
			\item \textbf{Simplifique:} \(3xy+4xy\)
			
			\pause
			\[
			3xy+4xy=7xy
			\]
			
			\pause
			\item \textbf{Evalúe:} \(5x+1\), si \(x=3\)
			
			\pause
			\[
			5(3)+1=15+1=16
			\]
			
			\pause
			\item \textbf{Patrón:} \(x,3x,9x,27x\). \pause Cada término se obtiene multiplicando por \(3\).
			
			\pause
			\item \textbf{Inconsistencia:} \(2x,4x,6x,11x\). \pause
			El término \(11x\) rompe el patrón de aumento de \(2x\) en \(2x\).
		\end{enumerate}
	\end{frame}
	
	% ---------------------------------
	\begin{frame}{Ejercicios en agronomía}
		\small
		\begin{enumerate}
			\item Simplifique:
			\[
			5x+3x-2x+x
			\]
			
			\pause
			\[
			5x+3x=8x,\quad 8x-2x=6x,\quad 6x+x=7x
			\]
			
			\pause
			\item Simplifique:
			\[
			4ab+6ab-3ab
			\]
			
			\pause
			\[
			4ab+6ab=10ab,\quad 10ab-3ab=7ab
			\]
			
			\pause
			\item Evalúe:
			\[
			2x+4 \quad \text{si } x=6
			\]
			
			\pause
			\[
			2(6)+4=12+4=16
			\]
			
			\pause
			\item Detecte patrón:
			\[
			2x,4x,8x,16x
			\]
			
			\pause
			Cada término se obtiene multiplicando por \(2\).
			
			\pause
			\item Detecte inconsistencia:
			\[
			3x,6x,9x,13x,15x
			\]
			
			\pause
			La inconsistencia es \(13x\), porque el patrón esperado aumenta de \(3x\) en \(3x\).
		\end{enumerate}
	\end{frame}
	
	\section{Patrones e inconsistencias}
	% ---------------------------------
	\begin{frame}{Detección de patrones}
		Observa la secuencia:
		\[
		2x,\ 4x,\ 6x,\ 8x
		\]
		
		\pause
		\textbf{¿Qué patrón identificas?}
		\begin{itemize}
			\item Aumenta de \(2x\) en \(2x\).
			\pause
			\item Existe una relación lineal.
			\pause
			\item Todos los términos son múltiplos de \(2x\).
		\end{itemize}
	\end{frame}
	
	% ---------------------------------
	\begin{frame}{Patrón agronómico}
		\Agronomia{
			Un registro de dosis de fertilizante por semana muestra:
			\[
			3x,\ 6x,\ 9x,\ 12x,\ 15x
			\]
		}
		
		\pause
		\textbf{Análisis:}
		\begin{itemize}
			\item Cada término aumenta en \(3x\).
			\pause
			\item La secuencia es lineal.
			\pause
			\item Todos los términos son múltiplos de \(3x\).
		\end{itemize}
		
		\pause
		\IdeaClave{Detectar patrones permite anticipar comportamientos y validar registros.}
	\end{frame}
	
	% ---------------------------------
	\begin{frame}{Detección de inconsistencias}
		Datos de riego:
		\[
		2x,\ 4x,\ 7x,\ 8x
		\]
		
		\pause
		\textbf{¿Qué ocurre aquí?}
		\begin{itemize}
			\item El término \(7x\) rompe el patrón esperado.
			\pause
			\item Podría existir un error de medición o de registro.
			\pause
			\item También podría indicar una condición particular no informada.
		\end{itemize}
	\end{frame}
	
	% ---------------------------------
	\begin{frame}{Inconsistencia agronómica}
		\Agronomia{
			Un técnico registra el volumen de riego aplicado en cinco sectores:
			\[
			5x,\ 10x,\ 15x,\ 19x,\ 25x
			\]
		}
		
		\pause
		\textbf{Patrón esperado:}
		\[
		5x,\ 10x,\ 15x,\ 20x,\ 25x
		\]
		
		\pause
		\textbf{Conclusión:}
		\begin{itemize}
			\item El término inconsistente es \(19x\).
			\pause
			\item Rompe el patrón de aumento regular en \(5x\).
			\pause
			\item Puede deberse a un error de transcripción o a una aplicación irregular.
		\end{itemize}
	\end{frame}
	
	\section{Análisis}
	% ---------------------------------
	\begin{frame}{Actividad de análisis}
		Situación:
		\[
		3x+2x+5x
		\]
		
		\pause
		\textbf{Pregunta:} ¿Es correcto decir que esto es igual a \(10x\)?
		
		\pause
		\textbf{Respuesta:} Sí.
		
		\pause
		Porque todos los términos son semejantes: tienen la misma variable \(x\) y el mismo exponente.
	\end{frame}
	
	% ---------------------------------
	\begin{frame}{Actividad de análisis ampliada}
		\Agronomia{
			Un técnico propone que
			\[
			4x^2+3x^2-2x
			\]
			se puede simplificar como
			\[
			5x^2.
			\]
		}
		
		\pause
		\textbf{¿Es correcto?}
		
		\pause
		No.
		
		\pause
		\textbf{Razón:}
		\begin{itemize}
			\item \(4x^2\) y \(3x^2\) sí son semejantes.
			\pause
			\item \(-2x\) no es semejante a \(x^2\), porque cambia el exponente.
		\end{itemize}
		
		\pause
		\textbf{Simplificación correcta:}
		\[
		4x^2+3x^2-2x=7x^2-2x
		\]
		
		\pause
		\IdeaClave{No se pueden mezclar términos con distinta parte literal.}
	\end{frame}
	
	\section{Aplicación integrada}
	% ---------------------------------
	\begin{frame}{Aplicación agronómica integrada}
		\begin{block}{Situación}
			En un cultivo experimental:
			\begin{itemize}
				\item se aplican \(2x\) litros de solución nutritiva en la mañana,
				\pause
				\item se aplican \(3x\) litros en la tarde,
				\pause
				\item se pierden \(x\) litros por filtración.
			\end{itemize}
		\end{block}
		
		\pause
		\textbf{Expresión algebraica:}
		\[
		2x+3x-x
		\]
		
		\pause
		\textbf{Simplificación:}
		\[
		4x
		\]
		
		\pause
		\IdeaClave{El aporte neto total al cultivo es de \(4x\) litros.}
	\end{frame}
	
	% ---------------------------------
	\begin{frame}{Aplicación integrada ampliada}
		\small
		\Agronomia{
			En una unidad productiva de pimientos:
			\begin{itemize}
				\item se aplican \(3x\) litros de fertilizante líquido en la mañana,
				\pause
				\item \(2x\) litros en la tarde,
				\pause
				\item se pierden \(x\) litros por drenaje,
				\pause
				\item se recuperan \(2x\) litros equivalentes mediante ajuste de fertirriego,
				\pause
				\item y existe una pérdida fija de 5 litros por mantenimiento.
			\end{itemize}
		}
		
		\pause
		\textbf{Expresión:}
		\[
		3x+2x-x+2x-5
		\]
		
		\pause
		\[
		5x-x+2x-5
		\]
		
		\pause
		\[
		4x+2x-5
		\]
		
		\pause
		\[
		6x-5
		\]
		
		\pause
		\IdeaClave{El balance final del sistema es \(6x-5\).}
	\end{frame}
	
	\section{Síntesis}
	% ---------------------------------
	\begin{frame}{Síntesis}
		\begin{itemize}[<+->]
			\item Un término algebraico incluye coeficiente, variable y exponente.
			\item Una expresión algebraica combina uno o más términos.
			\item Los términos semejantes se pueden sumar o restar directamente.
			\item Detectar patrones ayuda a interpretar datos de cultivo.
			\item Detectar inconsistencias permite prevenir errores en registros y decisiones agronómicas.
		\end{itemize}
	\end{frame}
	
	\section{Cierre}
	% ---------------------------------
	\begin{frame}{Cierre}
		\textbf{Hoy aprendimos:}
		\begin{itemize}[<+->]
			\item Qué es un término algebraico.
			\item Qué es una expresión algebraica.
			\item Cómo identificar términos semejantes.
			\item Cómo simplificar expresiones algebraicas.
			\item Cómo detectar patrones e inconsistencias.
		\end{itemize}
		
		\pause
		\begin{block}{Aplicación real}
			El álgebra permite modelar y optimizar recursos en sistemas agrícolas.
		\end{block}
	\end{frame}
	
	% ---------------------------------
	\begin{frame}{Fin de la clase}
		\centering
		\Huge ¿Preguntas?
	\end{frame}
	
\end{document}